%!TEX root = ../dissertation.tex

\chapter{Introduction}
\label{chp:intro}

Predictive process monitoring is a research discipline within the field of process mining that aims to provide stakeholders with a prediction on the outcome of an ongoing process \cite{vanderaalst2012}. These predictions are often generated through the analysis of data organized into event logs, which represent process executions as a sequence of events. Traditionally, these logs would represent the event data as inputs where each event is associated with exactly one object, or case identifier. While this representation was useful for modelling certain processes (such as insurance claims) which typically rely on a single object type being tied to each execution flow, other types of processes could not properly be described through this methodology without a significant loss of information. Object-centric processes, on the other hand, follow a different paradigm: rather than being executed in isolation, a process instance interacts with other instances of the same or different processes \cite{galanti2023a}. This paradigm aligns more closely with real-world applications, where a process unfolds through the parallel execution of different execution flows with different associated objects that interact with one another.
Contemporary predictive process monitoring research has begun focusing on performing their predictions by using Object-centric Event Logs (OCELs). OCELs represent process flows through the relationship between events and the multiple objects associated with each event step, as opposed to single-case-ID logs that force an artificial single-type view \cite{galanti2023a}. 

The type of process instances that relate to each other via many-to-many associations pose a significant challenge for predictive process monitoring. Since traditional process mining models assume that events are associated with exactly one object, these techniques cannot be directly applied to object-centric event data. Modern applications have leveraged deep learning models, such as neural networks, to address this issue. While traditional machine learning models such as Gradient Boosted Trees (GBT) or neural network architectures such as Long Short-Term Memory (LSTM) can successfully be applied to OCELs, they require the object-centric event data to be encoded in a tabular format which may lead to the loss or misrepresentation of information. Graph Neural Networks (GNNs) are specialized frameworks that extend traditional neural network methodologies so they may be used on graph-based datasets. By extending traditional neural networks to a graph-based domain these models are able to provide predictions based not only on the input's features but also on the underlying structure of the information, such as the ways events and objects interact throughout the execution of a process. 

Recent research has shown success in the application of graph-based encodings that respect the OCEL's underlying structure \cite{adams2022} and allow for GNNs to be used for predictive analysis. While these graph-based encodings originally used homogeneous graph structures, more recent works have explored whether heterogeneous structures can provide better predictive results. However, the evidence so far is inconclusive, with results showing improvement on some datasets but tying or underperforming on others \cite{smit2024} \cite{deleonieVolpato}. Similarly, while graph neural networks have proven to be extremely effective at providing accurate predictions on a variety of datasets their use in predictive process monitoring tasks has led to conflicting results, with some works citing a noticeable improvement \cite{adams2022} and others citing worse performance than traditional neural network models \cite{galantiDeleoni2024}. Additionally, graph neural networks (like their non-graph counterparts) offer limited transparency regarding the way they compute their predictions. These difficulties can hinder adoption in a business-process setting where a wrong or unexplained prediction has real operational consequences.

The need for providing end users with comprehensive explanations on the predictions made by deep learning models has led to the rise of Explainable Artificial Intelligence (XAI). XAI methodologies rely on the application of methods such as Shapley Value analysis to provide insight into the importance a model places on the different elements present in the input data. These methods have been successfully applied in other fields, such as Natural Language Processing, and have also found adoption in traditional (non-object-centric) processes. However, this thesis shows a gap in the application of explainability methods to predictive process monitoring on object-centric processes. In order to address this issue this thesis proposes an explainability framework composed of four complementary elements, each covering a distinct question a stakeholder might ask: which elements matter, which raw features drive them, what minimal change would change the prediction, and how much to trust the explanation. The first is a perturbation-based explanation which uses a combination of a Leave-One-Out methodology \cite{limonroejurafsky2017} to identify important elements and a Shapley Value analysis \cite{castro2009} to properly quantify their impact on a model's prediction, providing a clear look at what elements mattered in the input. Second, a feature/gradient-based explanation that complements the first method with a finer-grained view of which raw features drive the important elements previously highlighted \cite{sundararajan2017}. Third, a counterfactual explanation which demonstrates what changes to the factual input would be needed to achieve a different output prediction \cite{huang2022}, complementing the previous factual explanations with a more actionable look at what input changes can lead to a different prediction. And, finally, a set of fidelity metrics to assess the validity of the provided explanations \cite{stevens2022}, complementing the past three layers by giving the stakeholder a measure of how much they can trust the given explanations. The goal of this framework is not only to provide explanations on the chosen models, but also to present this information in a comprehensive manner to process stakeholders who may have no formal understanding of predictive methods as a whole. A more detailed discussion on each chosen component and their importance for the framework is provided in the state-of-the-art review in Chapter 2.

To validate the proposed framework, this thesis presents a set of experiments conducted on two synthetic OCEL 2.0 datasets. The experiments seek to train four distinct predictive models, two for each dataset, aimed at predicting the time remaining until the occurrence of a given event in the process. With the models properly trained and validated the explainability framework is applied to all four in order to identify the key elements driving each prediction, as well as quantify their impact on the prediction. These values are then presented in the form of graphs and tables in order to provide a stakeholder with an easy to interpret output that would help them better understand the predictive model. 

Overall, the results show that the proposed framework is able to identify important elements in the model and show it in a way that can generate trust in the stakeholder. A clear limitation of the framework is the associated computational cost of quantifying each element's impact on the prediction, but it is expected that further research on explainability frameworks for heterogeneous graphs will address this limitation.

The remainder of the document is organized as follows. Chapter 2 provides a review of related works on predictive process monitoring on object-centric event logs as well as explainability methods. Chapter 3 introduces the theoretical background. Chapter 4 provides a rundown of the work done to train the predictive models that were used to test the proposed explainability framework. Chapter 5 shows the results of the processes outlined in Chapter 4. Chapter 6 summarizes the findings and limitations of the proposed framework.